\chapter{Providers}
\label{ch:providers}

Providers configure LLM backends. They are declared at the top level and referenced by agents.

\section{Provider Declaration}

\begin{lstlisting}
provider provider_name {
    field: value
    field: value
}
\end{lstlisting}

\subsection{Cloud Providers}

\begin{lstlisting}
provider anthropic {
    api_key: env("ANTHROPIC_API_KEY")
    model: "claude-sonnet-4-20250514"
}

provider openai {
    api_key: env("OPENAI_API_KEY")
    model: "gpt-4o"
}
\end{lstlisting}

\subsection{Local Providers}

\begin{lstlisting}
provider local {
    backend: "ollama"
    host: "localhost:11434"
    model: "llama3:8b"
}
\end{lstlisting}

\section{Provider Fields}

\begin{table}[h]
\centering
\begin{tabular}{llll}
\toprule
\textbf{Field} & \textbf{Type} & \textbf{Required} & \textbf{Description} \\
\midrule
\code{model} & \type{string} & Yes & Default model identifier \\
\code{api\_key} & \type{string} & No & API authentication key \\
\code{backend} & \type{string} & No & Backend type (\code{"ollama"}, \code{"llama.cpp"}) \\
\code{host} & \type{string} & No & Backend host address \\
\code{temperature} & \type{float} & No & Default temperature (0.0--2.0) \\
\code{max\_tokens} & \type{int} & No & Default max tokens per response \\
\bottomrule
\end{tabular}
\caption{Provider fields}
\end{table}

\section{Environment Variables}

The \code{env()} function reads environment variables at runtime:

\begin{lstlisting}
provider anthropic {
    api_key: env("ANTHROPIC_API_KEY")
    model: env("ANTHROPIC_MODEL") or "claude-sonnet-4-20250514"
}
\end{lstlisting}

\begin{warning}
Never hardcode API keys in source files. Always use \code{env()} to read them from the environment.
\end{warning}

\section{Referencing Providers}

Providers are referenced by name in agent declarations:

\begin{lstlisting}
agent Writer {
    provider: anthropic      // uses anthropic provider
    system: "You are a writer."
}

agent FastHelper {
    provider: local          // uses local Ollama provider
    system: "You are a fast helper."
}
\end{lstlisting}

\section{Multiple Providers}

A program can declare multiple providers. Agents choose which to use:

\begin{lstlisting}
provider anthropic {
    api_key: env("ANTHROPIC_API_KEY")
    model: "claude-sonnet-4-20250514"
}

provider local {
    backend: "ollama"
    host: "localhost:11434"
    model: "llama3:8b"
}

// Use expensive cloud model for important tasks
agent Analyst {
    provider: anthropic
    system: "You analyze data carefully."
}

// Use cheap local model for simple tasks
agent Classifier {
    provider: local
    system: "Classify the input into categories."
}
\end{lstlisting}

\section{MCP Providers}

The Model Context Protocol (MCP) enables agents to use tools from external processes. MCP providers spawn a subprocess (or connect via HTTP) and discover its tools automatically.

\subsection{Stdio Transport}

The most common MCP transport. The runtime spawns a subprocess and communicates via JSON-RPC over stdin/stdout:

\begin{lstlisting}
provider filesystem {
    transport: "mcp"
    command: "npx"
    args: ["-y", "@modelcontextprotocol/server-filesystem", "/tmp"]
}
\end{lstlisting}

\subsection{SSE Transport}

For remote MCP servers accessible over HTTP:

\begin{lstlisting}
provider remote_tools {
    transport: "sse"
    endpoint: "http://localhost:9000/sse"
}
\end{lstlisting}

\subsection{MCP Provider Fields}

\begin{table}[h]
\centering
\begin{tabular}{llll}
\toprule
\textbf{Field} & \textbf{Type} & \textbf{Required} & \textbf{Description} \\
\midrule
\code{transport} & \type{string} & Yes & \code{"mcp"} (stdio) or \code{"sse"} (HTTP) \\
\code{command} & \type{string} & Stdio only & Executable to spawn \\
\code{args} & \type{[string]} & No & Command arguments \\
\code{endpoint} & \type{string} & SSE only & Server URL \\
\code{env} & \type{map} & No & Environment variables for subprocess \\
\code{headers} & \type{map} & No & HTTP headers (SSE only) \\
\bottomrule
\end{tabular}
\caption{MCP provider fields}
\end{table}

\subsection{Using MCP in Agents}

Agents reference MCP providers with the \code{mcp} field. The runtime discovers tools from all listed MCP providers and makes them available alongside regular tools:

\begin{lstlisting}
agent Assistant {
    provider: azure_openai
    system: "You are a helpful assistant with file system access."
    tools: [greet]              // regular tools
    mcp: [filesystem]           // MCP tools (auto-discovered)
    memory: conversation(max_turns: 10)
}
\end{lstlisting}

\subsection{MCP Server}

Haira workflows can be exposed as MCP tools for other applications to consume:

\begin{lstlisting}
import "mcp"

workflow Summarize(text: string) -> { summary: string } {
    """Summarize the given text into key points."""
    summary, err = Summarizer.ask(text)
    if err != nil { return { summary: "Error" } }
    return { summary: summary }
}

fn main() {
    mcp_server = mcp.Server([Summarize])
    mcp_server.listen(9000)  // SSE on port 9000
}
\end{lstlisting}

The \code{mcp.Server} exposes each workflow as an MCP tool. Tool names, descriptions, and input schemas are derived from the workflow declaration. The server supports both SSE (\code{.listen(port)}) and stdio (\code{.serve()}) transports.

\section{Provider Grammar}

\begin{lstlisting}[language=EBNF]
provider_decl  = "provider" identifier "{" { provider_field } "}" ;
provider_field = identifier ":" expression ;
\end{lstlisting}
